\documentclass[../DoAn.tex]{subfiles}
\begin{document}

Chương này trình bày quá trình từ thiết kế kiến trúc đến triển khai sản phẩm. Mở đầu là lựa chọn kiến trúc và các sơ đồ thiết kế chi tiết. Tiếp đó là kết quả xây dựng cùng các thành phần kỹ thuật trọng yếu của kho mã nguồn. Khép lại chương là kết quả kiểm thử hộp đen các chức năng nghiệp vụ và hướng dẫn triển khai.

\section{Thiết kế kiến trúc}
\label{section:4.1}

\subsection{Lựa chọn kiến trúc phần mềm}
\label{subsection:4.1.1}

Hệ thống áp dụng \textbf{kiến trúc phân tầng (Layered Architecture)} \cite{fowler2002poeaa,bass2021architecture} gồm bốn tầng từ trên xuống: \textbf{tầng trình bày} (Presentation) $\rightarrow$ \textbf{tầng nghiệp vụ} (Business) $\rightarrow$ \textbf{tầng truy cập dữ liệu} (Persistence) $\rightarrow$ \textbf{tầng dữ liệu} (Database). Trong cài đặt cụ thể, tầng trình bày được tổ chức nội bộ thành Route, Middleware và Controller; tầng nghiệp vụ là các Service; tầng truy cập dữ liệu gồm Repository và Prisma Client; tầng dữ liệu là PostgreSQL. Mỗi tầng chỉ tương tác với tầng liền kề, không gọi nhảy tầng. Đây \textbf{không phải mô hình MVC thuần}: backend là API thuần, không có tầng View (giao diện do frontend Next.js đảm nhiệm như một tầng client tách biệt), đồng thời bổ sung tầng Service và Repository theo mẫu Controller--Service--Repository để tách bạch rõ logic nghiệp vụ và truy cập dữ liệu.

Lớp Repository được tách khỏi Service trong một lần tái cấu trúc lớn của dự án. Trước đó, Service gọi trực tiếp \texttt{prisma.danhHieuHangNam.findMany(\dots)}, ràng buộc logic nghiệp vụ vào Prisma Client. Sau khi tách, Service gọi qua \texttt{danhHieuHangNamRepository.findMany(\dots)}: ranh giới rõ ràng, có thể thay thế ORM trong tương lai mà không phải sửa logic.

Các kiến trúc khác đã được cân nhắc nhưng không được chọn:

\begin{itemize}
    \item \textbf{Pure MVC} thiếu khái niệm Repository nên Service vừa chứa logic nghiệp vụ vừa thực hiện truy vấn, vi phạm nguyên tắc đơn trách nhiệm với một dự án có 23 model và nhiều quan hệ phức tạp như PM QLKT.
    \item \textbf{Clean Architecture / Hexagonal} đầy đủ với Use Case Layer, Entity Layer, Adapter Layer quá nặng cho quy mô dự án; lớp Use Case sẽ trùng lặp gần như hoàn toàn với Service trong kiến trúc phân tầng.
    \item \textbf{Microservices} không phù hợp với môi trường mạng nội bộ một máy chủ; chia tách dịch vụ chỉ làm tăng độ phức tạp triển khai mà không mang lại lợi ích về khả năng chịu tải.
\end{itemize}

\begin{figure}[H]
    \centering
    \includegraphics[width=0.62\textwidth]{KienTrucPhanTang}
    \caption{Sơ đồ kiến trúc phân tầng (Layered Architecture) bốn tầng}
    \label{fig:kientruc-phantang}
\end{figure}

\subsection{Tổng quan hệ thống}
\label{subsection:4.1.2}

Sản phẩm được tổ chức thành ba thành phần độc lập triển khai:

\begin{itemize}
    \item \textbf{Frontend (FE-QLKT)}: ứng dụng Next.js 14 chạy bằng Node.js. Cung cấp giao diện cho cả bốn vai trò người dùng. Kết nối tới backend qua REST API (HTTP) và Socket.IO (WebSocket).
    \item \textbf{Backend (BE-QLKT)}: API server Express + TypeScript chạy bằng Node.js. Cung cấp REST endpoint, Socket.IO server, thực hiện toàn bộ logic nghiệp vụ và quản lý phiên đăng nhập.
    \item \textbf{Cơ sở dữ liệu}: PostgreSQL chạy trong tiến trình riêng. Backend kết nối qua Prisma Client, không cho phép FE truy cập trực tiếp.
\end{itemize}

Cả ba thành phần này được triển khai trên cùng một máy chủ Linux của Học viện và giao tiếp qua localhost. Bên ngoài, người dùng truy cập qua trình duyệt thông qua mạng LAN; FE phục vụ giao diện, còn trình duyệt gọi trực tiếp tới BE qua REST API và Socket.IO.

\subsection{Thiết kế tổng quan các gói}
\label{subsection:4.1.3}

Cấu trúc thư mục mã nguồn được tổ chức theo nguyên tắc đơn trách nhiệm và đặt tên theo tiếng Anh chuẩn để dễ tham chiếu trong các sơ đồ.

\subsubsection{Gói phía frontend}

Mã nguồn frontend được tổ chức bên trong thư mục \texttt{FE-QLKT/src/}. Gói \texttt{app/} chứa cây định tuyến của Next.js App Router với bốn nhánh tương ứng bốn vai trò người dùng (super-admin, admin, manager, user). Gói \texttt{components/} chứa các thành phần React dùng chung, được tách thành các thư mục con theo miền nghiệp vụ: \texttt{auth/} cho màn hình đăng nhập và đổi mật khẩu, \texttt{proposals/} cho biểu mẫu tạo đề xuất nhiều bước với mỗi loại trong bảy loại đề xuất, \texttt{personnel/} cho danh sách và chi tiết quân nhân, \texttt{system-logs/} cho bảng nhật ký hệ thống. Gói \texttt{contexts/} đặt \texttt{AuthContext} quản lý trạng thái đăng nhập trong toàn ứng dụng. Gói \texttt{hooks/} chứa các React hook tự viết như \texttt{useFetch}, \texttt{useAuthGuard}, \texttt{useSocket}. Gói \texttt{lib/} chứa các module phụ trợ: \texttt{api/} chia theo từng miền nghiệp vụ và đóng gói các yêu cầu REST tới backend, \texttt{award/} chứa các hàm trợ giúp render chuỗi danh hiệu phía frontend, \texttt{schemas.ts} định nghĩa các schema Zod xác thực biểu mẫu. Gói \texttt{constants/} chứa các hằng số dùng chung như danh sách vai trò, trạng thái, danh hiệu.

\begin{figure}[H]
    \centering
    \includegraphics[width=0.9\textwidth]{goi-fe}
    \caption{Sơ đồ gói phía frontend}
    \label{fig:goi-frontend}
\end{figure}

\subsubsection{Gói phía backend}

Mã nguồn backend được tổ chức bên trong thư mục \texttt{BE-QLKT/src/} theo các tầng kiến trúc đã trình bày ở mục \ref{subsection:4.1.1}. Gói \texttt{routes/} định nghĩa các điểm cuối REST API, mỗi miền nghiệp vụ một tệp. Gói \texttt{middlewares/} chứa các tệp \texttt{auth.ts} (xác thực và phân quyền), \texttt{auditLog.ts} (ghi nhật ký mọi thao tác làm thay đổi dữ liệu), \texttt{unitFilter.ts} (lọc dữ liệu theo phạm vi đơn vị của Manager) và \texttt{validate.ts} (bọc Zod validation). Gói \texttt{controllers/} chứa các tệp điều hướng yêu cầu HTTP từ route tới service, được giữ mỏng (mỗi phương thức dưới 50 dòng mã). Gói \texttt{services/} chứa logic nghiệp vụ chính, được tách tiếp thành các thư mục con: \texttt{proposal/strategies/} cho strategy registry của bảy loại đề xuất, \texttt{eligibility/} cho hàm \texttt{chainEligibility} xét quy tắc chuỗi cá nhân và đơn vị, \texttt{profile/annual.ts} cho logic chuỗi cá nhân và \texttt{profile/unit.ts} cho logic chuỗi đơn vị. Gói \texttt{repositories/} đóng gói toàn bộ truy cập Prisma, mỗi model một tệp \texttt{.repository.ts}. Gói \texttt{helpers/} chứa các tiện ích phụ trợ chia theo nhóm: \texttt{auditLog/}, \texttt{notification/}, \texttt{excel/}, \texttt{awardValidation/}. Gói \texttt{validations/} chứa các schema Zod cho từng route. Tệp \texttt{prisma/schema.prisma} ở cấp gốc của module backend định nghĩa toàn bộ 23 model.

\begin{figure}[H]
    \centering
    \includegraphics[width=0.9\textwidth]{goi-be}
    \caption{Sơ đồ gói phía backend}
    \label{fig:goi-backend}
\end{figure}

\subsection{Thiết kế chi tiết gói nghiệp vụ Khen thưởng cá nhân hằng năm}
\label{subsection:4.1.4}

Để minh hoạ cách các tầng phối hợp trong một module nghiệp vụ cụ thể, phần này phân tích module chuỗi danh hiệu hằng năm cá nhân, module phức tạp nhất hệ thống. Sơ đồ chi tiết các gói thành phần và quan hệ phụ thuộc của module được trình bày ở Hình~\ref{fig:goi-chi-tiet-hangnam}.

\begin{figure}[H]
    \centering
    \includegraphics[width=\textwidth]{so-do-goi-cnhn}
    \caption{Sơ đồ chi tiết gói nghiệp vụ Khen thưởng cá nhân hằng năm}
    \label{fig:goi-chi-tiet-hangnam}
\end{figure}

Khi người dùng truy cập trang hồ sơ một quân nhân, dòng dữ liệu đi qua các tầng theo thứ tự:

\begin{enumerate}
    \item \textbf{Route} \texttt{GET /api/profiles/annual/:personnel\_id} (định nghĩa tại \texttt{routes/profile.route.ts}) tiếp nhận yêu cầu, đi qua middleware \texttt{verifyToken}; quyền xem hồ sơ của chính mình được kiểm ở tầng service nên route không cần guard theo vai trò.
    \item \textbf{Controller} \texttt{profileController.getAnnualProfile} (tại \texttt{controllers/profile.controller.ts}) đọc tham số \texttt{personnel\_id}, gọi service tương ứng, gói kết quả qua \texttt{ResponseHelper.success}.
    \item \textbf{Service} \texttt{profileService.recalculateAnnualProfile} (tại \texttt{services/profile/annual.ts}) thực thi logic: nạp danh sách \texttt{DanhHieuHangNam} qua repository, gọi \texttt{computeChainContext} để dẫn xuất ngữ cảnh chuỗi, lặp qua \texttt{PERSONAL\_CHAIN\_AWARDS} và gọi \texttt{checkChainEligibility} cho mỗi cấp, ghi kết quả vào bảng suy diễn \texttt{HoSoHangNam}.
    \item \textbf{Eligibility module} \texttt{services/eligibility/chainEligibility.ts} chứa hàm \texttt{checkChainEligibility(config, context, hasReceivedLifetime)} là hàm thuần, không truy vấn cơ sở dữ liệu, dễ kiểm chứng với mọi bộ đầu vào.
    \item \textbf{Repository} \texttt{repositories/danhHieu.repository.ts} cung cấp các phương thức \texttt{findManyByPersonnelId}, \texttt{upsertByPersonnelYear}.
    \item \textbf{Prisma Client} thực thi truy vấn SQL tới bảng \texttt{danh\_hieu\_hang\_nam}.
\end{enumerate}

Việc tách Service (logic nghiệp vụ) khỏi Eligibility module (quy tắc thuần) là chủ ý: module quy tắc hoàn toàn thuần nên có thể kiểm chứng với mọi kịch bản đầu vào mà không cần kết nối cơ sở dữ liệu, còn Service chỉ điều phối (nạp dữ liệu, gọi quy tắc, ghi kết quả). Ranh giới rõ ràng, dễ bảo trì.

\subsection{Thiết kế chi tiết gói nghiệp vụ Khen thưởng đơn vị hằng năm}
\label{subsection:goi-chi-tiet-donvi}

Gói khen thưởng đơn vị hằng năm có cấu trúc tầng tương tự gói cá nhân, được trình bày ở Hình~\ref{fig:goi-chi-tiet-donvi}. Điểm khác biệt nằm ở phân quyền và dữ liệu nghiệp vụ, không nằm ở kiến trúc.

\begin{figure}[H]
    \centering
    \includegraphics[width=\textwidth]{so-do-goi-dvhn}
    \caption{Sơ đồ chi tiết gói nghiệp vụ Khen thưởng đơn vị hằng năm}
    \label{fig:goi-chi-tiet-donvi}
\end{figure}

Dòng dữ liệu khi tính lại điều kiện khen thưởng cho một đơn vị đi qua các tầng:

\begin{enumerate}
    \item \textbf{Route} (tại \texttt{routes/unitAnnualAward.route.ts}) đi qua middleware \texttt{verifyToken} và \texttt{requireAdminOrManager} --- khác với gói cá nhân, khen thưởng đơn vị chỉ Admin/Manager được thao tác nên route chặn vai trò ngay từ đầu.
    \item \textbf{Controller} \texttt{unitAnnualAwardController} chuyển yêu cầu tới service.
    \item \textbf{Service} \texttt{unitAnnualAwardService} (lớp facade) uỷ thác cho module \texttt{services/unitAnnualAward/eligibility.ts}: nạp danh sách \texttt{DanhHieuDonViHangNam} qua repository, đếm số năm Đơn vị Quyết thắng liên tục, rồi gọi quy tắc chuỗi và ghi kết quả vào bảng suy diễn \texttt{HoSoDonViHangNam}.
    \item \textbf{Eligibility module} \texttt{services/eligibility/unitChainEvaluator.ts} bọc hàm chung \texttt{checkChainEligibility} với bộ cấu hình \texttt{UNIT\_CHAIN\_AWARDS}.
    \item \textbf{Repository} \texttt{danhHieuDonViHangNamRepository} và \texttt{unitAnnualProfileRepository} đóng gói truy cập Prisma tới hai bảng \texttt{danh\_hieu\_don\_vi\_hang\_nam} và \texttt{ho\_so\_don\_vi\_hang\_nam}.
\end{enumerate}

Điểm thiết kế đáng chú ý là cả hai gói (cá nhân và đơn vị) cùng dùng chung hàm thuần \texttt{checkChainEligibility} ở tầng \texttt{eligibility/}, chỉ khác bộ cấu hình (\texttt{PERSONAL\_CHAIN\_AWARDS} so với \texttt{UNIT\_CHAIN\_AWARDS}). Nhờ vậy, quy tắc xét chuỗi danh hiệu được mô tả một chỗ duy nhất, hai luồng không thể lệch nhau, và việc bổ sung danh hiệu mới chỉ cần thêm cấu hình. Khác biệt nghiệp vụ thể hiện rõ trên sơ đồ: gói cá nhân có thêm bảng \texttt{ThanhTichKhoaHoc} (yêu cầu nghiên cứu khoa học mỗi năm) và lấy dữ liệu theo từng quân nhân, trong khi gói đơn vị không xét nghiên cứu khoa học.

\subsection{Thiết kế lớp}
\label{subsection:4.1.5}

Phần này trình bày sơ đồ lớp cho năm module nghiệp vụ chính. Mỗi sơ đồ chỉ ra các lớp/interface cốt lõi cùng các phương thức quan trọng và mối quan hệ.

\textbf{Module xử lý chuỗi danh hiệu.} Trung tâm là interface \texttt{ChainAwardConfig} mô tả cấu hình một cấp trong chuỗi: mã danh hiệu, số năm chu kỳ, danh sách cờ tiền điều kiện, yêu cầu NCKH, có phải danh hiệu một lần duy nhất hay không, tên cột cờ trong cơ sở dữ liệu, nhãn chuỗi. Hai mảng cấu hình \texttt{PERSONAL\_CHAIN\_AWARDS} và \texttt{UNIT\_CHAIN\_AWARDS} chứa các cấp tương ứng cho cá nhân và đơn vị. Hàm \texttt{checkChainEligibility(config, context, hasReceivedLifetime, flagsInWindow)} đọc cấu hình, đối chiếu với ngữ cảnh thực tế và trả về \texttt{EligibilityResult \{ eligible: boolean, reason: string \}}. Interface \texttt{ChainContext} đóng gói các chỉ số dẫn xuất từ lịch sử: độ dài chuỗi, năm gần nhất nhận từng cấp, số cờ tiền điều kiện trong khung xét.

\textbf{Module Strategy đề xuất.} Interface \texttt{ProposalStrategy} áp dụng mẫu thiết kế Strategy \cite{gof1994patterns}, định nghĩa bốn phương thức bắt buộc cho mỗi loại đề xuất: \texttt{buildSubmitPayload}, \texttt{validateApprove}, \texttt{importInTransaction}, \texttt{buildSuccessMessage}. Bảy lớp triển khai (\texttt{caNhanHangNamStrategy}, \texttt{donViHangNamStrategy}, \texttt{hccsvvStrategy}, \texttt{hcqkqtStrategy}, \texttt{kncStrategy}, \texttt{hcbvtqStrategy}, \texttt{nckhStrategy}) đăng ký vào \texttt{REGISTRY} trong \texttt{services/proposal/strategies/index.ts}. Hai lớp single-medal (\texttt{hcqkqtStrategy} và \texttt{kncStrategy}) chia sẻ logic qua hàm trợ giúp \texttt{singleMedalImporter} để tránh trùng lặp.

\textbf{Module xác thực và phân quyền.} Lớp chính là \texttt{AuthService} cung cấp các phương thức \texttt{login}, \texttt{refresh}, \texttt{logout}, \texttt{changePassword}. Middleware \texttt{verifyToken} đọc header \texttt{Authorization: Bearer}, xác minh Access Token, gắn \texttt{req.user}. Middleware \texttt{requireRole(...allowed)} factory trả về middleware kiểm tra \texttt{req.user.role} thuộc danh sách cho phép. Middleware \texttt{unitFilter} truy vấn \texttt{co\_quan\_don\_vi\_id}/\texttt{don\_vi\_truc\_thuoc\_id} của người dùng và áp giới hạn phạm vi cho mỗi yêu cầu của Manager.

\textbf{Module quản lý đề xuất.} \texttt{ProposalService} điều phối luồng tạo và phê duyệt đề xuất, dùng \texttt{proposalRepository} để truy vấn và cập nhật. Hàm \texttt{approveProposal} tại \texttt{services/proposal/approve.ts} đảm nhiệm logic phê duyệt phức tạp; nó gọi tuần tự các module phụ \texttt{approve/validation.ts} (tiền kiểm), \texttt{approve/decisionMappings.ts} (gắn quyết định) và \texttt{approve/import.ts} (ghi dữ liệu) bên trong cùng một giao dịch (transaction) Prisma.

\textbf{Module nhập và xuất Excel.} Bao gồm các hàm trợ giúp dùng chung (\texttt{loadWorkbook}, \texttt{getAndValidateWorksheet}, \texttt{parseHeaderMap}, \texttt{sanitizeRowData}) và các strategy chuyên biệt cho từng loại nghiệp vụ. Mỗi strategy có hai phương thức \texttt{previewImport} (chỉ kiểm tra) và \texttt{confirmImport} (mở giao dịch ghi).

\begin{figure}[H]
    \centering
    \includegraphics[width=\textwidth]{lop-to-chuc}
    \caption{Sơ đồ lớp module Tổ chức, Quân nhân và Tài khoản}
    \label{fig:lop-modules}
\end{figure}

\begin{figure}[H]
    \centering
    \includegraphics[width=\textwidth]{lop-de-xuat}
    \caption{Sơ đồ lớp module Đề xuất khen thưởng (mẫu Strategy)}
    \label{fig:lop-dexuat}
\end{figure}

\begin{figure}[H]
    \centering
    \includegraphics[width=\textwidth]{lop-khen-thuong}
    \caption{Sơ đồ lớp module Khen thưởng, Hồ sơ điều kiện và Quyết định}
    \label{fig:lop-khenthuong}
\end{figure}

\subsection{Biểu đồ tuần tự}
\label{subsection:4.1.6}

Hai luồng nghiệp vụ phức tạp và tiêu biểu nhất là tạo đề xuất và phê duyệt khen thưởng. Cả hai được mô hình hoá bằng sơ đồ tuần tự (sequence diagram), thể hiện tương tác giữa các đối tượng theo thời gian.

\textbf{Luồng tạo đề xuất} (Hình~\ref{fig:tuantu-dexuat}): Manager hoàn thành biểu mẫu năm bước $\rightarrow$ FE gọi \texttt{POST /api/proposals} $\rightarrow$ chuỗi middleware $\rightarrow$ \texttt{proposalService.submitProposal} $\rightarrow$ \texttt{getProposalStrategy(type).buildSubmitPayload} $\rightarrow$ \texttt{proposalRepository.create} $\rightarrow$ phát sự kiện Socket.IO tới Admin $\rightarrow$ trả về đề xuất đã tạo.

\begin{figure}[H]
    \centering
    \includegraphics[width=0.95\textwidth]{seq-de-xuat}
    \caption{Sơ đồ tuần tự luồng Đề xuất khen thưởng}
    \label{fig:tuantu-dexuat}
\end{figure}

\textbf{Luồng phê duyệt đề xuất} (Hình~\ref{fig:tuantu-pheduyet}): Admin mở chi tiết đề xuất $\rightarrow$ bấm ``Phê duyệt'' $\rightarrow$ FE gọi \texttt{POST /api/proposals/:id/approve} $\rightarrow$ chuỗi middleware $\rightarrow$ \texttt{approveProposal} mở giao dịch Prisma $\rightarrow$ \texttt{validation.preflight} $\rightarrow$ \texttt{strategy.validateApprove} $\rightarrow$ \texttt{decisionMappings.attach} $\rightarrow$ \texttt{strategy.importInTransaction} $\rightarrow$ cập nhật \texttt{BangDeXuat.status = APPROVED} $\rightarrow$ tính lại hồ sơ suy diễn $\rightarrow$ commit $\rightarrow$ ghi nhật ký hệ thống $\rightarrow$ phát sự kiện Socket.IO.

\begin{figure}[H]
    \centering
    \includegraphics[width=0.95\textwidth]{seq-phe-duyet}
    \caption{Sơ đồ tuần tự luồng Phê duyệt khen thưởng}
    \label{fig:tuantu-pheduyet}
\end{figure}

\subsection{Thiết kế cơ sở dữ liệu}
\label{subsection:4.1.7}

Cơ sở dữ liệu được mô hình hoá thành 23 model trong tệp \texttt{prisma/schema.prisma}. Hình \ref{fig:erd-tongthe} là sơ đồ ERD tổng thể của hệ thống.

\begin{figure}[H]
    \centering
    \includegraphics[width=\textwidth,height=0.85\textheight,keepaspectratio]{DB}
    \caption{Sơ đồ ERD tổng thể của 23 model}
    \label{fig:erd-tongthe}
\end{figure}

Sơ đồ ERD tổng thể bao quát toàn bộ 23 bảng với các quan hệ khoá ngoại chính. Có thể nhận diện ba ``trục'' chính: (1) trục tổ chức gồm \texttt{CoQuanDonVi}, \texttt{DonViTrucThuoc}, \texttt{ChucVu}, \texttt{LichSuChucVu}; (2) trục quân nhân gồm \texttt{QuanNhan}, \texttt{TaiKhoan}; (3) trục khen thưởng gồm \texttt{BangDeXuat}, \texttt{DanhHieuHangNam}, \texttt{KhenThuongHCCSVV}, \texttt{HuanChuongQuanKyQuyetThang}, \texttt{KyNiemChuongVSNXDQDNDVN}, \texttt{KhenThuongHCBVTQ}, \texttt{KhenThuongDotXuat}, \texttt{ThanhTichKhoaHoc}. Các bảng suy diễn \texttt{HoSoHangNam}, \texttt{HoSoNienHan}, \texttt{HoSoCongHien}, \texttt{HoSoDonViHangNam} chứa kết quả tính toán mà hệ thống tự động cập nhật lại sau mỗi thay đổi nguồn.

Sáu bảng mô tả chi tiết schema sau đây được chọn vì là các bảng trọng yếu nhất, sử dụng nhiều cấu trúc dữ liệu phức tạp (JSON, mảng, kiểu thời gian).

\begin{table}[H]
\centering
\caption{Schema bảng \texttt{QuanNhan} (Personnel)}
\label{tab:schema-quannhan}
\small
\renewcommand{\arraystretch}{1.2}
\begin{tabular}{|p{3.8cm}|p{2.3cm}|p{3.5cm}|p{3.4cm}|}
\hline
\textbf{Cột} & \textbf{Kiểu} & \textbf{Ràng buộc} & \textbf{Mô tả} \\ \hline
\texttt{id} & String & PK, CUID & Khoá chính \\ \hline
\texttt{cccd} & VarChar(20) & UNIQUE, NULL & Số căn cước công dân \\ \hline
\dbid{ho_ten} & VarChar(255) & NOT NULL & Họ và tên đầy đủ \\ \hline
\dbid{ngay_sinh} & Date & NULL & Ngày sinh \\ \hline
\dbid{gioi_tinh} & VarChar(10) & NULL & Giới tính (NAM, NU) \\ \hline
\dbid{ngay_nhap_ngu} & Date & NULL & Ngày nhập ngũ \\ \hline
\dbid{ngay_xuat_ngu} & Date & NULL & Ngày xuất ngũ (nếu đã rời) \\ \hline
\dbid{cap_bac} & VarChar(50) & NULL & Cấp bậc hiện tại \\ \hline
\dbid{co_quan_don_vi_id} & VarChar(30) & FK $\rightarrow$ CoQuanDonVi & Liên kết tới cơ quan đơn vị \\ \hline
\dbid{don_vi_truc_thuoc_id} & VarChar(30) & FK $\rightarrow$ DonViTrucThuoc & Liên kết tới đơn vị trực thuộc \\ \hline
\dbid{chuc_vu_id} & VarChar(30) & FK $\rightarrow$ ChucVu & Chức vụ hiện tại \\ \hline
\texttt{createdAt} & Timestamp & DEFAULT now() & Thời điểm tạo \\ \hline
\texttt{updatedAt} & Timestamp & AUTO UPDATE & Thời điểm cập nhật \\ \hline
\end{tabular}
\end{table}

Ràng buộc nghiệp vụ: chỉ một trong \texttt{co\_quan\_don\_vi\_id} và \texttt{don\_vi\_truc\_thuoc\_id} được điền. Khi cập nhật chuyển đơn vị, cần dùng giao dịch để cập nhật cả bộ đếm \texttt{so\_luong} của hai đơn vị liên quan.

\begin{table}[H]
\centering
\caption{Schema bảng \texttt{BangDeXuat} (Proposal)}
\label{tab:schema-dexuat}
\small
\renewcommand{\arraystretch}{1.2}
\begin{tabular}{|p{3.8cm}|p{2.3cm}|p{3.5cm}|p{3.4cm}|}
\hline
\textbf{Cột} & \textbf{Kiểu} & \textbf{Ràng buộc} & \textbf{Mô tả} \\ \hline
\texttt{id} & String & PK, CUID & \\ \hline
\dbid{loai_de_xuat} & VarChar(20) & NOT NULL & Một trong 7 mã loại đề xuất \\ \hline
\dbid{nguoi_de_xuat_id} & VarChar(30) & FK $\rightarrow$ TaiKhoan & Người tạo đề xuất \\ \hline
\texttt{nam} & Int & NOT NULL & Năm xét \\ \hline
\texttt{thang} & Int & NULL & Tháng xét (chỉ với một số loại) \\ \hline
\texttt{status} & VarChar(20) & DEFAULT 'PENDING' & PENDING, APPROVED, REJECTED \\ \hline
\dbid{data_danh_hieu} & JSON & NULL & Dữ liệu chi tiết hằng năm \\ \hline
\dbid{data_thanh_tich} & JSON & NULL & Dữ liệu thành tích NCKH \\ \hline
\dbid{data_nien_han} & JSON & NULL & Dữ liệu HCCSVV \\ \hline
\dbid{data_cong_hien} & JSON & NULL & Dữ liệu HCBVTQ \\ \hline
\dbid{files_attached} & JSON & NULL & Mảng file đính kèm \\ \hline
\dbid{rejection_reason} & Text & NULL & Lý do từ chối (nếu có) \\ \hline
\dbid{nguoi_duyet_id} & VarChar(30) & FK $\rightarrow$ TaiKhoan, NULL & Admin đã phê duyệt \\ \hline
\dbid{ngay_duyet} & Timestamp & NULL & Ngày phê duyệt \\ \hline
\end{tabular}
\end{table}

Việc dùng JSON cho 4 trường \texttt{data\_*} là chủ ý. Mỗi loại đề xuất có lược đồ con riêng nên việc tạo nhiều bảng trung gian sẽ phá vỡ tính đồng nhất của bảng \texttt{BangDeXuat}. JSON đủ để đảm bảo lược đồ linh hoạt; xác thực được đảm nhiệm ở tầng Zod và Strategy tương ứng từng loại.

\begin{table}[H]
\centering
\caption{Schema bảng \texttt{DanhHieuHangNam} (Annual Award)}
\label{tab:schema-dhhn}
\small
\renewcommand{\arraystretch}{1.2}
\begin{tabular}{|p{3.8cm}|p{2.3cm}|p{3.5cm}|p{3.4cm}|}
\hline
\textbf{Cột} & \textbf{Kiểu} & \textbf{Ràng buộc} & \textbf{Mô tả} \\ \hline
\texttt{id} & String & PK, CUID & \\ \hline
\dbid{quan_nhan_id} & VarChar(30) & FK $\rightarrow$ QuanNhan & \\ \hline
\texttt{nam} & Int & UNIQUE (quan\_nhan\_id, nam) & Năm xét danh hiệu \\ \hline
\dbid{danh_hieu} & VarChar(20) & NULL & Mã danh hiệu cơ bản (CSTĐCS, CSTT) \\ \hline
\dbid{nhan_bkbqp} & Boolean & DEFAULT false & Đã nhận BKBTBQP trong năm \\ \hline
\dbid{so_quyet_dinh_bkbqp} & VarChar(100) & NULL & Số quyết định BKBTBQP \\ \hline
\dbid{nhan_cstdtq} & Boolean & DEFAULT false & Đã nhận CSTĐTQ trong năm \\ \hline
\dbid{so_quyet_dinh_cstdtq} & VarChar(100) & NULL & Số quyết định CSTĐTQ \\ \hline
\dbid{nhan_bkttcp} & Boolean & DEFAULT false & Đã nhận BKTTCP trong năm \\ \hline
\dbid{so_quyet_dinh_bkttcp} & VarChar(100) & NULL & Số quyết định BKTTCP \\ \hline
\dbid{ghi_chu} & Text & NULL & Ghi chú bổ sung \\ \hline
\end{tabular}
\end{table}

Ràng buộc duy nhất \texttt{(quan\_nhan\_id, nam)} đảm bảo mỗi quân nhân chỉ có một bản ghi cho một năm. Các trường \texttt{nhan\_*} cho phép cộng dồn các cấp danh hiệu chuỗi vào cùng năm.

\begin{table}[H]
\centering
\caption{Schema bảng \texttt{LichSuChucVu} (Position History)}
\label{tab:schema-lsc}
\small
\renewcommand{\arraystretch}{1.2}
\begin{tabular}{|p{3.8cm}|p{2.3cm}|p{3.5cm}|p{3.4cm}|}
\hline
\textbf{Cột} & \textbf{Kiểu} & \textbf{Ràng buộc} & \textbf{Mô tả} \\ \hline
\texttt{id} & String & PK, CUID & \\ \hline
\dbid{quan_nhan_id} & VarChar(30) & FK $\rightarrow$ QuanNhan & \\ \hline
\dbid{chuc_vu_id} & VarChar(30) & FK $\rightarrow$ ChucVu & \\ \hline
\dbid{ngay_bat_dau} & Date & NOT NULL & Ngày bắt đầu giữ chức vụ \\ \hline
\dbid{ngay_ket_thuc} & Date & NULL & NULL = đang giữ \\ \hline
\dbid{he_so_chuc_vu} & Float & DEFAULT 0 & Hệ số chức vụ (snapshot) \\ \hline
\dbid{so_thang} & Int & NULL & Số tháng giữ chức vụ \\ \hline
\end{tabular}
\end{table}

Hàm tính lại số tháng tự động cập nhật trường \texttt{so\_thang} mỗi khi có thay đổi mốc thời gian, phục vụ thuật toán xét HCBVTQ.

\begin{table}[H]
\centering
\caption{Schema bảng \texttt{SystemLog} (Audit Log)}
\label{tab:schema-syslog}
\small
\renewcommand{\arraystretch}{1.2}
\begin{tabular}{|p{3.8cm}|p{2.3cm}|p{3.5cm}|p{3.4cm}|}
\hline
\textbf{Cột} & \textbf{Kiểu} & \textbf{Ràng buộc} & \textbf{Mô tả} \\ \hline
\texttt{id} & String & PK, CUID & \\ \hline
\dbid{nguoi_thuc_hien_id} & VarChar(30) & FK $\rightarrow$ TaiKhoan, NULL & NULL với hành động hệ thống \\ \hline
\dbid{actor_role} & VarChar(20) & NOT NULL & Vai trò tại thời điểm thao tác \\ \hline
\texttt{action} & VarChar(50) & NOT NULL & CREATE, UPDATE, DELETE, ... \\ \hline
\texttt{resource} & VarChar(50) & NOT NULL & Loại tài nguyên \\ \hline
\dbid{tai_nguyen_id} & VarChar(30) & NULL & Mã định danh tài nguyên \\ \hline
\texttt{description} & VarChar(500) & NOT NULL & Mô tả tiếng Việt \\ \hline
\texttt{payload} & JSON & NULL & Dữ liệu trước, sau khi thay đổi \\ \hline
\dbid{ip_address} & VarChar(45) & NULL & Địa chỉ IPv4/IPv6 \\ \hline
\dbid{user_agent} & Text & NULL & Trình duyệt người dùng \\ \hline
\texttt{createdAt} & Timestamp & DEFAULT now() & Thời điểm ghi log \\ \hline
\end{tabular}
\end{table}

\begin{table}[H]
\centering
\caption{Schema bảng \texttt{HoSoHangNam} (Annual Profile, Derived)}
\label{tab:schema-hsnn}
\small
\renewcommand{\arraystretch}{1.2}
\begin{tabular}{|p{3.8cm}|p{2.3cm}|p{3.5cm}|p{3.4cm}|}
\hline
\textbf{Cột} & \textbf{Kiểu} & \textbf{Ràng buộc} & \textbf{Mô tả} \\ \hline
\texttt{id} & String & PK, CUID & \\ \hline
\dbid{quan_nhan_id} & VarChar(30) & UNIQUE, FK $\rightarrow$ QuanNhan & Mỗi quân nhân một bản ghi \\ \hline
\dbid{cstdcs_lien_tuc} & Int & DEFAULT 0 & Độ dài chuỗi CSTĐCS hiện tại \\ \hline
\dbid{nckh_lien_tuc} & Int & DEFAULT 0 & Số năm NCKH liên tục \\ \hline
\dbid{bkbqp_lien_tuc} & Int & DEFAULT 0 & Năm liên tục từ lần BKBTBQP cuối \\ \hline
\dbid{cstdtq_lien_tuc} & Int & DEFAULT 0 & Năm liên tục từ lần CSTĐTQ cuối \\ \hline
\dbid{du_dieu_kien_bkbqp} & Boolean & DEFAULT false & Đủ điều kiện đề nghị BKBTBQP \\ \hline
\dbid{du_dieu_kien_cstdtq} & Boolean & DEFAULT false & Đủ điều kiện đề nghị CSTĐTQ \\ \hline
\dbid{du_dieu_kien_bkttcp} & Boolean & DEFAULT false & Đủ điều kiện đề nghị BKTTCP \\ \hline
\dbid{goi_y} & Text & NULL & Thông điệp gợi ý tiếng Việt \\ \hline
\texttt{updatedAt} & Timestamp & UPDATED & Thời điểm tính lại gần nhất \\ \hline
\end{tabular}
\end{table}

Toàn văn schema (618 dòng, 23 model) được tổ chức trong tệp \texttt{BE-QLKT/prisma/schema.prisma} của kho mã nguồn dự án.

\section{Thiết kế chi tiết}
\label{section:4.2}

\subsection{Thiết kế giao diện}
\label{subsection:4.2.1}

Giao diện hệ thống được phác thảo bằng wireframe trước khi triển khai bằng React + Ant Design. Bảy wireframe dưới đây thể hiện bố cục các trang trọng tâm, từ màn hình đăng nhập đến các trang nghiệp vụ chính.

\textbf{Màn hình đăng nhập} (Hình~\ref{fig:wf-login}) bố trí một thẻ đăng nhập căn giữa trên nền ảnh, gồm logo, tiêu đề, hai trường tên đăng nhập và mật khẩu cùng nút đăng nhập. Đây là điểm vào duy nhất của hệ thống. Mọi vai trò đều xác thực qua màn hình này trước khi được điều hướng tới trang chủ tương ứng.

\begin{figure}[H]
    \centering
    \includegraphics[width=0.7\textwidth]{wireframe/wireframe-login}
    \caption{Wireframe màn hình đăng nhập}
    \label{fig:wf-login}
\end{figure}

\textbf{Trang chủ Admin} (Hình~\ref{fig:wf-dashboard}) có cấu trúc ba khu vực: thanh điều hướng bên trái (cây cơ quan đơn vị và đơn vị trực thuộc có thể thu gọn), khu vực dashboard ở giữa với bốn thẻ thống kê (tổng quân nhân, tổng đơn vị, đề xuất chờ duyệt, khen thưởng tháng này) cùng các biểu đồ phân bố, khu vực phụ bên phải với danh sách thông báo thời gian thực. Bố cục ba khu vực này được giữ nhất quán xuyên suốt các trang nghiệp vụ chính (Quản lý quân nhân, Quản lý đề xuất, Quản lý quyết định) để giảm thời gian làm quen cho người dùng mới.

\begin{figure}[H]
    \centering
    \includegraphics[width=0.85\textwidth]{wireframe/wireframe-dashboard}
    \caption{Wireframe trang chủ Admin với bốn thẻ thống kê và các biểu đồ}
    \label{fig:wf-dashboard}
\end{figure}

\textbf{Trang danh sách quân nhân} (Hình~\ref{fig:wf-personnel-list}) gồm thanh điều hướng bên trái, các thẻ thống kê nhanh, thanh tìm kiếm và bộ lọc nhiều tiêu chí ở đầu, cùng bảng dữ liệu phân trang có cột thao tác cuối dòng. Mẫu bố cục ``bộ lọc trên, bảng dưới, thao tác cuối dòng'' này được tái sử dụng cho hầu hết các trang danh sách (quân nhân, đề xuất, quyết định, khen thưởng).

\begin{figure}[H]
    \centering
    \includegraphics[width=0.85\textwidth]{wireframe/wireframe-personnel-list}
    \caption{Wireframe trang danh sách quân nhân}
    \label{fig:wf-personnel-list}
\end{figure}

\textbf{Trang chi tiết quân nhân} (Hình~\ref{fig:wf-personnel-detail}) hiển thị thông tin định danh ở đầu (ảnh đại diện, họ tên, mã và vai trò) và chia phần thân thành các tab tương ứng các nhóm khen thưởng: Hằng năm, Niên hạn, Cống hiến, Quân kỳ, Kỷ niệm chương, NCKH, Đột xuất. Mỗi tab có bảng dữ liệu kèm thanh tiến trình hiển thị các chu kỳ đang đến mốc đề nghị (ví dụ ``BKBTBQP: 1/2 năm CSTĐCS đã tích luỹ''). Cách tiếp cận trực quan này giúp Manager nhanh chóng nhận biết quân nhân nào sắp đủ điều kiện.

\begin{figure}[H]
    \centering
    \includegraphics[width=0.85\textwidth]{wireframe/wireframe-personnel-detail}
    \caption{Wireframe trang chi tiết quân nhân với các tab khen thưởng}
    \label{fig:wf-personnel-detail}
\end{figure}

\textbf{Trang danh sách đề xuất} (Hình~\ref{fig:wf-proposal-list}) bổ sung hàng tab trạng thái (Tất cả, Chờ duyệt, Đã duyệt, Từ chối) phía trên bảng, mỗi dòng kèm nhãn trạng thái màu và nút xem chi tiết. Nút ``Tạo mới'' ở góc phải dẫn tới biểu mẫu tạo đề xuất.

\begin{figure}[H]
    \centering
    \includegraphics[width=0.85\textwidth]{wireframe/wireframe-proposal-list}
    \caption{Wireframe trang danh sách đề xuất khen thưởng}
    \label{fig:wf-proposal-list}
\end{figure}

\textbf{Biểu mẫu tạo đề xuất} (Hình~\ref{fig:wf-proposal-create}) áp dụng biểu mẫu năm bước (multi-step wizard) với thanh stepper ở đầu. Bước 1: chọn loại đề xuất (radio button có icon đại diện cho từng nhóm) và năm/tháng. Bước 2: chọn quân nhân hoặc đơn vị từ bảng có lọc theo nhiều tiêu chí. Bước 3: thiết lập danh hiệu cụ thể cho từng đối tượng đã chọn, kèm gợi ý từ hệ thống. Bước 4: tải lên tệp đính kèm minh chứng. Bước 5: xem lại toàn bộ thông tin và gửi đề xuất. Người dùng có thể quay lại các bước trước mà không mất dữ liệu nhờ trạng thái được lưu vào \texttt{sessionStorage} của trình duyệt.

\begin{figure}[H]
    \centering
    \includegraphics[width=0.85\textwidth]{wireframe/wireframe-proposal-create}
    \caption{Wireframe biểu mẫu tạo đề xuất năm bước}
    \label{fig:wf-proposal-create}
\end{figure}

\textbf{Trang danh sách khen thưởng} (Hình~\ref{fig:wf-awards}) dùng hàng tab chọn loại danh hiệu phía trên, kèm thanh lọc và bảng dữ liệu bên dưới. Cùng một bố cục phục vụ cả bảy loại khen thưởng, chỉ thay đổi các cột hiển thị theo loại đang chọn.

\begin{figure}[H]
    \centering
    \includegraphics[width=0.85\textwidth]{wireframe/wireframe-awards}
    \caption{Wireframe trang danh sách khen thưởng theo loại}
    \label{fig:wf-awards}
\end{figure}

Ngoài bảy trang trên, trang nhập Excel (vùng kéo thả tệp kèm bảng xem trước hai tab ``Hợp lệ'' / ``Lỗi'', tab Lỗi tô đỏ và liệt kê chỉ số dòng) và trang nhật ký hệ thống (bảng nhóm theo ngày, mỗi dòng mở rộng xem \texttt{payload} JSON chi tiết) đều kế thừa cùng khung bố cục ba khu vực nói trên.

\subsection{Hệ thống thiết kế (Design System)}
\label{subsection:4.2.2}

Bảng màu của hệ thống kế thừa mặc định của Ant Design 5 (gam xanh đại dương cho hành động chính, đỏ cho hành động xoá, vàng cho cảnh báo, xanh lá cho thành công) với một số tinh chỉnh nhỏ qua Tailwind CSS để phù hợp đặc thù môi trường quân sự (giảm độ bão hoà của các màu phụ để giao diện trông trang trọng hơn).

Hệ chữ chính dùng phông \texttt{Inter} (sans-serif) cho phần thân nội dung do khả năng hiển thị tốt với tiếng Việt có dấu, và phông \texttt{Source Code Pro} cho khu vực hiển thị mã quyết định và mã định danh kỹ thuật (CCCD, mã quyết định) để dễ đọc các chuỗi số dài.

Khoảng cách (spacing) được chuẩn hoá theo thang 4 px của Tailwind. Tất cả các thẻ (Card) có lề trong tối thiểu 16 px để tránh nội dung sát mép. Bán kính bo góc được chọn 6 px cho mọi phần tử để tạo cảm giác hiện đại nhưng vẫn nghiêm túc.

\section{Xây dựng ứng dụng}
\label{section:4.3}

\subsection{Thư viện và công cụ sử dụng}
\label{subsection:4.3.1}

Sản phẩm dùng các thư viện mã nguồn mở liệt kê tại Bảng \ref{tab:thuvien}, phân theo ba nhóm: tầng giao diện, tầng máy chủ và các tiện ích vận hành.

\begin{table}[H]
\centering
\caption{Thư viện và công cụ chính}
\label{tab:thuvien}
\begin{tabular}{|c|p{6cm}|p{4.5cm}|c|}
\hline
\textbf{STT} & \textbf{Mục đích} & \textbf{Thư viện / Công cụ} & \textbf{Phiên bản} \\ \hline
1 & Khung phát triển giao diện người dùng & Next.js & 14.x \\ \hline
2 & Thư viện thành phần UI & Ant Design & 5.x \\ \hline
3 & Khung CSS utility & Tailwind CSS & 3.x \\ \hline
4 & Khung phát triển phía máy chủ & Express & 4.x \\ \hline
5 & Ngôn ngữ kiểu tĩnh & TypeScript & 5.x \\ \hline
6 & Hệ quản trị CSDL & PostgreSQL & 15 \\ \hline
7 & ORM & Prisma & 6.x \\ \hline
8 & Xác thực JWT & jsonwebtoken & 9.x \\ \hline
9 & Băm mật khẩu & bcrypt & 5.x \\ \hline
10 & Xác thực dữ liệu (BE + FE) & Zod & 4.x \\ \hline
11 & Thông báo thời gian thực & Socket.IO & 4.x \\ \hline
12 & Đọc / ghi Excel & ExcelJS & 4.x \\ \hline
14 & Lập lịch nội bộ & node-cron & 4.x \\ \hline
15 & Quản lý tiến trình production & PM2 & 5.x \\ \hline
16 & Header bảo mật HTTP & Helmet & 8.x \\ \hline
17 & CORS & cors & 2.x \\ \hline
18 & Giới hạn tốc độ & express-rate-limit & 8.x \\ \hline
19 & Quản lý tệp tải lên & Multer & 1.x \\ \hline
\end{tabular}
\end{table}

Môi trường phát triển dùng Visual Studio Code với extension Prisma, ESLint và Prettier; mã nguồn quản lý qua Git và GitHub; cơ sở dữ liệu được tra cứu trực quan qua Prisma Studio (\texttt{npx prisma studio}).

\subsection{Kết quả đạt được}
\label{subsection:4.3.2}

Sau quá trình xây dựng, hệ thống đã hoàn thiện các chức năng chính sau:

\begin{itemize}
    \item Quản lý quân nhân: thêm, sửa, xoá thủ công qua biểu mẫu (không nhập/xuất Excel); chuyển đơn vị và đổi chức vụ chỉ Admin, giữ nguyên lịch sử khen thưởng; phân quyền theo phạm vi đơn vị cho Manager.
    \item Quản lý đơn vị (cơ quan đơn vị và đơn vị trực thuộc) và chức vụ: cây tổ chức hai cấp; tự cập nhật bộ đếm \texttt{so\_luong} khi quân nhân chuyển đơn vị.
    \item Quản lý lịch sử chức vụ: nhập danh sách giai đoạn giữ chức vụ với hệ số; tự tính \texttt{so\_thang} mỗi khi thay đổi mốc thời gian.
    \item Đề xuất khen thưởng cho cả bảy loại với biểu mẫu năm bước; tự động gợi ý quân nhân đủ điều kiện theo từng loại.
    \item Phê duyệt đề xuất với giao dịch, gắn số quyết định, tải lên PDF, ghi nhật ký đầy đủ trạng thái trước và sau khi thay đổi.
    \item Tính lại điều kiện khen thưởng hằng năm cá nhân và đơn vị tự động sau mỗi thay đổi nguồn; gợi ý đề nghị bằng tiếng Việt.
    \item Nhập Excel hai bước (xem trước và xác nhận với giao dịch) cho cả bảy loại nghiệp vụ; xuất danh sách khen thưởng theo nhiều tiêu chí.
    \item Quản trị tài khoản bốn cấp; đặt lại mật khẩu; nhật ký hệ thống; sao lưu định kỳ; khu vực DevZone cho SuperAdmin.
    \item Thông báo thời gian thực qua Socket.IO khi có đề xuất mới, được phê duyệt, bị từ chối, bị xoá hoặc nhập Excel hoàn tất.
\end{itemize}

Tổng quy mô mã nguồn TypeScript khoảng 96.000 dòng (không tính dòng trống và bình luận), chia thành 38.000 dòng phía backend và 58.000 dòng phía frontend, cộng với 618 dòng \texttt{schema.prisma}. Kho mã nguồn quản lý theo monorepo với hai thư mục \texttt{BE-QLKT/} và \texttt{FE-QLKT/} ở mức gốc.

\subsection{Minh hoạ các chức năng chính}
\label{subsection:4.3.3}

Phần này trình bày các ảnh chụp màn hình tiêu biểu của sản phẩm trong môi trường phát triển. Toàn bộ ảnh được chụp trên trình duyệt Chromium ở độ phân giải 1920×1080, sau đó cắt khu vực có nội dung và lưu định dạng PNG. Các hình được nhóm theo bốn vai trò người dùng (Quản trị viên, Phòng Chính trị, Chỉ huy đơn vị, Người dùng) cùng nhóm màn hình chung; mỗi vai trò có trang chủ (dashboard) riêng với bộ thẻ thống kê và biểu đồ phù hợp phạm vi dữ liệu của vai trò đó.

\subsubsection{Màn hình chung và xác thực}

\textbf{Trang chủ} (Hình~\ref{fig:ui-home}) là trang giới thiệu công khai, nêu tên hệ thống cùng vài số liệu tổng quan và hai lối vào ``Đăng nhập'' / ``Đề xuất khen thưởng''. \textbf{Màn hình đăng nhập} (Hình~\ref{fig:ui-login}) là điểm vào duy nhất của hệ thống: thẻ đăng nhập kính mờ căn giữa trên nền ảnh Học viện, gồm logo, trường tài khoản và mật khẩu, liên kết quên mật khẩu cùng nút đăng nhập. Mọi vai trò xác thực qua đây trước khi được điều hướng tới trang chủ tương ứng.

\begin{figure}[H]
    \centering
    \includegraphics[width=0.85\textwidth]{UI/01-trang-chu}
    \caption{Trang chủ giới thiệu hệ thống}
    \label{fig:ui-home}
\end{figure}

\begin{figure}[H]
    \centering
    \includegraphics[width=0.85\textwidth]{UI/02-dang-nhap}
    \caption{Màn hình đăng nhập}
    \label{fig:ui-login}
\end{figure}

\subsubsection{Vai trò Quản trị viên (SuperAdmin)}

\textbf{Trang chủ Quản trị viên} (Hình~\ref{fig:ui-sa-dashboard}) tổng hợp toàn hệ thống với bốn thẻ thống kê (tổng tài khoản, tài khoản mới trong 30 ngày, hoạt động trong 7 ngày, tổng nhật ký) cùng các biểu đồ phân bố vai trò, hoạt động hệ thống và top hành động phổ biến.

\begin{figure}[H]
    \centering
    \includegraphics[width=0.92\textwidth]{UI/03-superadmin-dashboard}
    \caption{Trang chủ vai trò Quản trị viên}
    \label{fig:ui-sa-dashboard}
\end{figure}

\textbf{Quản lý tài khoản} (Hình~\ref{fig:ui-sa-account}) liệt kê toàn bộ tài khoản kèm vai trò, quân nhân liên kết, đơn vị và thao tác sửa, xoá. \textbf{Tạo tài khoản} (Hình~\ref{fig:ui-sa-account-create}) là biểu mẫu chọn vai trò, đặt tên đăng nhập (gợi ý dùng CCCD), mật khẩu (để trống sẽ dùng mật khẩu mặc định) và gắn cơ quan, đơn vị, chức vụ.

\begin{figure}[H]
    \centering
    \includegraphics[width=0.92\textwidth]{UI/04-superadmin-quan-ly-tai-khoan}
    \caption{Quản lý danh sách tài khoản}
    \label{fig:ui-sa-account}
\end{figure}

\begin{figure}[H]
    \centering
    \includegraphics[width=0.78\textwidth]{UI/05-superadmin-tao-tai-khoan}
    \caption{Biểu mẫu tạo tài khoản mới}
    \label{fig:ui-sa-account-create}
\end{figure}

\textbf{Quản lý cơ quan đơn vị} (Hình~\ref{fig:ui-sa-unit}) hiển thị cây tổ chức hai cấp (cơ quan đơn vị và các đơn vị trực thuộc); mỗi đơn vị có thể mở rộng để xem danh sách quân nhân trực thuộc kèm cấp bậc và chức vụ.

\begin{figure}[H]
    \centering
    \includegraphics[width=0.92\textwidth]{UI/06-superadmin-quan-ly-co-quan-don-vi}
    \caption{Quản lý cơ quan đơn vị và chức vụ}
    \label{fig:ui-sa-unit}
\end{figure}

\subsubsection{Vai trò Phòng Chính trị (Admin)}

\textbf{Trang chủ Phòng Chính trị} (Hình~\ref{fig:ui-admin-dashboard}) trình bày năm thẻ thống kê phạm vi toàn quân số quản lý (tổng quân nhân, đơn vị, chức vụ, đề xuất chờ duyệt, đề xuất bị từ chối) cùng các biểu đồ NCKH và đề xuất theo loại và trạng thái.

\begin{figure}[H]
    \centering
    \includegraphics[width=0.92\textwidth]{UI/27-admin-dashboard}
    \caption{Trang chủ vai trò Phòng Chính trị}
    \label{fig:ui-admin-dashboard}
\end{figure}

\textbf{Chi tiết quân nhân} chia phần thân thành nhiều tab. Tab \textit{Thông tin cơ bản} (Hình~\ref{fig:ui-admin-personnel-info}) hiển thị thông tin định danh, thông tin cá nhân và thông tin Đảng. Tab \textit{Hồ sơ khen thưởng} (Hình~\ref{fig:ui-admin-personnel-award}) trình bày hồ sơ niên hạn (HCCSVV ba hạng) và cống hiến (HCBVTQ) kèm thanh tiến trình và gợi ý điều kiện cho từng cấp danh hiệu.

\begin{figure}[H]
    \centering
    \includegraphics[width=0.92\textwidth]{UI/07-admin-chi-tiet-quan-nhan-thong-tin}
    \caption{Chi tiết quân nhân -- tab Thông tin cơ bản}
    \label{fig:ui-admin-personnel-info}
\end{figure}

\begin{figure}[H]
    \centering
    \includegraphics[width=0.92\textwidth]{UI/08-admin-chi-tiet-quan-nhan-ho-so-khen-thuong}
    \caption{Chi tiết quân nhân -- tab Hồ sơ khen thưởng với thanh tiến trình điều kiện}
    \label{fig:ui-admin-personnel-award}
\end{figure}

\textbf{Duyệt đề xuất} gồm trang danh sách (Hình~\ref{fig:ui-admin-proposal-list}) với tab trạng thái và bộ lọc năm, loại, cùng cửa sổ chi tiết (Hình~\ref{fig:ui-admin-proposal-detail}) hiển thị tệp đính kèm cùng bảng danh hiệu của đề xuất kèm số quyết định trước khi phê duyệt hoặc từ chối.

\begin{figure}[H]
    \centering
    \includegraphics[width=0.92\textwidth]{UI/09-admin-duyet-de-xuat-danh-sach}
    \caption{Danh sách đề xuất chờ phê duyệt}
    \label{fig:ui-admin-proposal-list}
\end{figure}

\begin{figure}[H]
    \centering
    \includegraphics[width=0.92\textwidth]{UI/10-admin-duyet-de-xuat-chi-tiet}
    \caption{Chi tiết một đề xuất khen thưởng}
    \label{fig:ui-admin-proposal-detail}
\end{figure}

\textbf{Quản lý quyết định} gồm bảng danh sách quyết định (Hình~\ref{fig:ui-admin-decision-list}) với số QĐ, năm, ngày ký, người ký và thao tác (xem, sửa, xoá, xem PDF), cùng cửa sổ chi tiết (Hình~\ref{fig:ui-admin-decision-detail}) hiển thị đầy đủ thông tin và nút xem file PDF.

\begin{figure}[H]
    \centering
    \includegraphics[width=0.92\textwidth]{UI/11-admin-quan-ly-quyet-dinh}
    \caption{Quản lý danh sách quyết định khen thưởng}
    \label{fig:ui-admin-decision-list}
\end{figure}

\begin{figure}[H]
    \centering
    \includegraphics[width=0.85\textwidth]{UI/12-admin-chi-tiet-quyet-dinh}
    \caption{Chi tiết một quyết định khen thưởng}
    \label{fig:ui-admin-decision-detail}
\end{figure}

\textbf{Quản lý khen thưởng} (Hình~\ref{fig:ui-admin-award-list}) dùng hàng tab chọn loại danh hiệu, bộ lọc năm, họ tên, danh hiệu và bảng dữ liệu kèm số quyết định; cửa sổ \textbf{xuất Excel} (Hình~\ref{fig:ui-admin-award-export}) cho phép chọn khoảng thời gian, đơn vị, quân nhân và danh hiệu cần xuất.

\begin{figure}[H]
    \centering
    \includegraphics[width=0.92\textwidth]{UI/13-admin-quan-ly-khen-thuong}
    \caption{Quản lý khen thưởng theo loại danh hiệu}
    \label{fig:ui-admin-award-list}
\end{figure}

\begin{figure}[H]
    \centering
    \includegraphics[width=0.85\textwidth]{UI/14-admin-xuat-excel-khen-thuong}
    \caption{Cửa sổ xuất Excel danh sách khen thưởng}
    \label{fig:ui-admin-award-export}
\end{figure}

\textbf{Thêm khen thưởng đồng loạt} (Hình~\ref{fig:ui-admin-bulk}) áp dụng biểu mẫu năm bước với stepper, cho phép tải file mẫu Excel rồi tải lên danh sách để nhập hàng loạt. \textbf{Khen thưởng đột xuất} (Hình~\ref{fig:ui-admin-adhoc}) quản lý các khen thưởng ngoài chu kỳ theo năm, đối tượng, hình thức và số quyết định.

\begin{figure}[H]
    \centering
    \includegraphics[width=0.92\textwidth]{UI/15-admin-them-khen-thuong-dong-loat}
    \caption{Thêm khen thưởng đồng loạt qua Excel}
    \label{fig:ui-admin-bulk}
\end{figure}

\begin{figure}[H]
    \centering
    \includegraphics[width=0.92\textwidth]{UI/16-admin-khen-thuong-dot-xuat}
    \caption{Quản lý khen thưởng đột xuất}
    \label{fig:ui-admin-adhoc}
\end{figure}

\textbf{Nhật ký hệ thống} (Hình~\ref{fig:ui-admin-logs}) gồm panel bộ lọc (chip lọc nhanh theo thời gian, ô tìm kiếm, khoảng ngày, vai trò, hành động), bốn thẻ thống kê số lượng theo loại hành động và bảng nhật ký có thể sắp xếp theo cột, mỗi dòng mô tả chi tiết thao tác đã thực hiện.

\begin{figure}[H]
    \centering
    \includegraphics[width=0.92\textwidth]{UI/17-admin-nhat-ky-he-thong}
    \caption{Trang nhật ký hệ thống với bộ lọc và thống kê}
    \label{fig:ui-admin-logs}
\end{figure}

\subsubsection{Vai trò Chỉ huy đơn vị (Manager)}

\textbf{Trang chủ Chỉ huy đơn vị} (Hình~\ref{fig:ui-manager-dashboard}) giới hạn dữ liệu trong phạm vi đơn vị quản lý, với năm thẻ thống kê (quân số, tổng CSTĐCS, tổng NCKH, khen thưởng, đề xuất chờ duyệt) và các biểu đồ khen thưởng, đề xuất, NCKH.

\begin{figure}[H]
    \centering
    \includegraphics[width=0.92\textwidth]{UI/28-manager-dashboard}
    \caption{Trang chủ vai trò Chỉ huy đơn vị}
    \label{fig:ui-manager-dashboard}
\end{figure}

\textbf{Đề xuất khen thưởng} (Hình~\ref{fig:ui-manager-proposal}) liệt kê các đề xuất đã gửi kèm hộp hướng dẫn, tab trạng thái và nút tạo đề xuất mới. \textbf{Biểu mẫu tạo đề xuất} áp dụng wizard năm bước: bước 1 chọn loại khen thưởng từ các thẻ bảy loại (Hình~\ref{fig:ui-manager-proposal-b1}); bước 2 chọn quân nhân từ bảng có lọc và cột hiển thị điều kiện (Hình~\ref{fig:ui-manager-proposal-b2}).

\begin{figure}[H]
    \centering
    \includegraphics[width=0.92\textwidth]{UI/18-manager-de-xuat-khen-thuong}
    \caption{Danh sách đề xuất khen thưởng của đơn vị}
    \label{fig:ui-manager-proposal}
\end{figure}

\begin{figure}[H]
    \centering
    \includegraphics[width=0.92\textwidth]{UI/19-manager-tao-de-xuat-b1-chon-loai}
    \caption{Tạo đề xuất -- Bước 1: chọn loại khen thưởng}
    \label{fig:ui-manager-proposal-b1}
\end{figure}

\begin{figure}[H]
    \centering
    \includegraphics[width=0.92\textwidth]{UI/20-manager-tao-de-xuat-b2-chon-quan-nhan}
    \caption{Tạo đề xuất -- Bước 2: chọn quân nhân}
    \label{fig:ui-manager-proposal-b2}
\end{figure}

\textbf{Lịch sử chức vụ} (Hình~\ref{fig:ui-manager-position}) ghi nhận các giai đoạn giữ chức vụ kèm hệ số và ba thẻ tổng thời gian theo nhóm hệ số. \textbf{Danh hiệu hằng năm} của một quân nhân (Hình~\ref{fig:ui-manager-annual}) liệt kê các danh hiệu theo từng năm kèm số quyết định, còn \textbf{thành tích NCKH} (Hình~\ref{fig:ui-manager-nckh}) liệt kê các đề tài/sáng kiến theo năm.

\begin{figure}[H]
    \centering
    \includegraphics[width=0.92\textwidth]{UI/21-manager-lich-su-chuc-vu}
    \caption{Lịch sử chức vụ của quân nhân}
    \label{fig:ui-manager-position}
\end{figure}

\begin{figure}[H]
    \centering
    \includegraphics[width=0.92\textwidth]{UI/22-manager-khen-thuong-hang-nam}
    \caption{Danh hiệu khen thưởng hằng năm của quân nhân}
    \label{fig:ui-manager-annual}
\end{figure}

\begin{figure}[H]
    \centering
    \includegraphics[width=0.92\textwidth]{UI/23-manager-thanh-tich-nckh}
    \caption{Thành tích nghiên cứu khoa học của quân nhân}
    \label{fig:ui-manager-nckh}
\end{figure}

\textbf{Chi tiết quân nhân -- tab Quản lý chi tiết} (Hình~\ref{fig:ui-manager-personnel-manage}) tập hợp các lối tắt quản lý theo lưới: lịch sử chức vụ, danh hiệu hằng năm, NCKH và các huân chương, huy chương, kỷ niệm chương.

\begin{figure}[H]
    \centering
    \includegraphics[width=0.92\textwidth]{UI/24-manager-chi-tiet-quan-nhan-quan-ly}
    \caption{Chi tiết quân nhân -- tab Quản lý chi tiết}
    \label{fig:ui-manager-personnel-manage}
\end{figure}

\subsubsection{Vai trò Người dùng (User)}

\textbf{Trang chủ Người dùng} (Hình~\ref{fig:ui-user-dashboard}) hiển thị hồ sơ cá nhân của quân nhân đăng nhập: các thẻ thống kê cá nhân (tổng CSTĐCS, NCKH, CSTĐCS liên tục, số tháng cống hiến), thông tin cá nhân và tóm tắt hồ sơ điều kiện hằng năm và niên hạn kèm gợi ý. \textbf{Hồ sơ danh hiệu cá nhân hằng năm} (Hình~\ref{fig:ui-user-annual}) trình bày thanh tiến trình tổng hợp và dòng thời gian chuỗi danh hiệu theo năm (CSTĐCS $\rightarrow$ BKBTBQP $\rightarrow$ CSTĐTQ $\rightarrow$ BKTTCP).

\begin{figure}[H]
    \centering
    \includegraphics[width=0.92\textwidth]{UI/25-user-dashboard}
    \caption{Trang chủ vai trò Người dùng}
    \label{fig:ui-user-dashboard}
\end{figure}

\begin{figure}[H]
    \centering
    \includegraphics[width=0.92\textwidth]{UI/26-user-ho-so-danh-hieu-hang-nam}
    \caption{Hồ sơ danh hiệu cá nhân hằng năm với dòng thời gian chuỗi danh hiệu}
    \label{fig:ui-user-annual}
\end{figure}

\subsection{Tóm tắt các thành phần kỹ thuật trọng yếu}
\label{subsection:4.3.4}

Phần này tóm tắt năm thành phần kỹ thuật trọng yếu của kho mã nguồn để người đọc nắm cách tổ chức bên trong, không trích nguyên đoạn mã.

\textbf{Hàm \texttt{computeChainContext} dẫn xuất ngữ cảnh chuỗi.} Đặt tại \texttt{BE-QLKT/src/services/profile/annual.ts}. Đầu vào gồm danh sách bản ghi \texttt{DanhHieuHangNam} của một quân nhân, độ dài chuỗi CSTĐCS hiện tại và năm xét. Đầu ra là đối tượng \texttt{ChainContext} chứa năm bắt đầu chuỗi, năm gần nhất nhận BKBTBQP / CSTĐTQ / BKTTCP còn nằm trong cửa sổ chuỗi, số năm CSTĐCS tích luỹ từ lần cuối nhận mỗi cấp và số chu kỳ đã lỡ. Hàm hoàn toàn thuần (không truy vấn cơ sở dữ liệu) nên dễ kiểm chứng với nhiều bộ dữ liệu mẫu phủ các kịch bản từ chu kỳ vừa đủ đến lỡ nhiều chu kỳ.

\textbf{Chuỗi middleware bảo mật năm tầng.} Mỗi yêu cầu làm thay đổi dữ liệu đi qua tuần tự: \texttt{verifyToken} (xác thực Access Token), \texttt{checkRole} (phân quyền vai trò), \texttt{writeLimiter} (giới hạn tốc độ), middleware tải lên tệp đính kèm (\texttt{multer}), \texttt{auditLog} (ghi nhật ký), rồi mới đến controller. Mỗi tầng giải quyết một mối quan tâm độc lập và tái sử dụng được giữa các route, phản ánh nguyên tắc đơn trách nhiệm.

\textbf{Repository tách Prisma Client.} Mỗi repository (ví dụ \texttt{danhHieuRepository}) đóng gói thao tác Prisma sau các phương thức nghiệp vụ (\texttt{findByPersonnelId}, \texttt{upsertByPersonnelYear}, \texttt{deleteManyByPersonnelId}). Phương thức nhận tham số tuỳ chọn \texttt{tx} kiểu \texttt{PrismaLike = prisma | Prisma.TransactionClient}, cho phép tham gia vào giao dịch do Service mở. Service có thể \texttt{prisma.\$transaction(async tx => \dots)} rồi truyền \texttt{tx} qua nhiều lời gọi repository, đảm bảo tính nguyên tử (all-or-nothing) của giao dịch.

\textbf{Hàm \texttt{checkChainEligibility} thực thi quy tắc chuỗi.} Đặt tại \texttt{BE-QLKT/src/services/eligibility/chainEligibility.ts}. Đầu vào gồm cấu hình cấp danh hiệu (\texttt{ChainAwardConfig}), ngữ cảnh chuỗi (\texttt{ChainStreaks}), cờ ``đã nhận'' (\texttt{hasReceived}) và số cờ tiền điều kiện trong khung xét (\texttt{FlagsInWindow}). Logic ba bước: chặn nhận lại nếu là BKTTCP cá nhân đã có; kiểm tra đủ chu kỳ (\texttt{streakLength} $\ge$ \texttt{cycleYears} và là bội số của \texttt{cycleYears}); kiểm tra đủ cờ tiền điều kiện trong khung xét cùng yêu cầu NCKH (nếu cấp danh hiệu yêu cầu). Trả về \texttt{EligibilityResult \{ eligible, reason \}} kèm thông điệp tiếng Việt.

\textbf{Strategy registry cho bảy loại đề xuất.} Tại \texttt{BE-QLKT/src/services/proposal/strategies/index.ts}, đối tượng \texttt{REGISTRY} ánh xạ mỗi \texttt{ProposalType} sang một thực thể tuân giao diện \texttt{ProposalStrategy}. Bảy strategy đã đăng ký gồm \texttt{caNhanHangNamStrategy}, \texttt{donViHangNamStrategy}, \texttt{hccsvvStrategy}, \texttt{hcqkqtStrategy}, \texttt{kncStrategy}, \texttt{hcbvtqStrategy}, \texttt{nckhStrategy}; hai strategy đơn huân chương (\texttt{hcqkqtStrategy}, \texttt{kncStrategy}) chia sẻ logic qua \texttt{singleMedalImporter}. Hàm \texttt{requireProposalStrategy(type)} trả về strategy tương ứng hoặc ném lỗi nếu chưa đăng ký. Để bổ sung loại đề xuất mới chỉ cần tạo tệp triển khai mới, thêm một dòng đăng ký vào \texttt{REGISTRY} và bổ sung khoá vào hằng số \texttt{PROPOSAL\_TYPES}.

\section{Kiểm thử}
\label{section:4.4}

\subsection{Kiểm thử hộp đen các chức năng nghiệp vụ}
\label{subsection:4.4.1}

Các chức năng nghiệp vụ chính được kiểm thử hộp đen thủ công thông qua giao diện trên môi trường phát triển. Tám bảng kiểm thử dưới đây tóm tắt kết quả.

\begin{table}[H]
\centering
\caption{Kiểm thử chức năng Đăng nhập}
\label{tab:kt-dangnhap}
\begin{tabular}{|c|p{3.8cm}|p{3.3cm}|p{3.8cm}|c|}
\hline
\textbf{STT} & \textbf{Chức năng} & \textbf{Đầu vào} & \textbf{Đầu ra mong muốn} & \textbf{Kết quả} \\ \hline
1 & Đăng nhập đúng & username + password đúng & Chuyển hướng tới trang chủ Admin & Đạt \\ \hline
2 & Sai mật khẩu & username đúng + password sai & Thông báo chung không tiết lộ trường nào sai & Đạt \\ \hline
3 & Tài khoản không tồn tại & username = noexist & Thông báo chung không lộ thông tin & Đạt \\ \hline
4 & Vượt rate limit \texttt{authLimiter} & Vượt 30 yêu cầu thất bại / IP / 5 phút & Trả mã 429, từ chối tạm thời & Đạt \\ \hline
5 & Refresh token & Access Token hết hạn & Tự động xin token mới & Đạt \\ \hline
\end{tabular}
\end{table}

\begin{table}[H]
\centering
\caption{Kiểm thử chức năng Quản lý quân nhân}
\label{tab:kt-quannhan}
\begin{tabular}{|c|p{3.8cm}|p{3.3cm}|p{3.8cm}|c|}
\hline
\textbf{STT} & \textbf{Chức năng} & \textbf{Đầu vào} & \textbf{Đầu ra mong muốn} & \textbf{Kết quả} \\ \hline
1 & Thêm mới hợp lệ & Đầy đủ trường & Bản ghi mới có trong danh sách & Đạt \\ \hline
2 & Trùng CCCD & CCCD đã tồn tại & Hiển thị lỗi ``Số CCCD đã tồn tại'' & Đạt \\ \hline
3 & Thiếu trường bắt buộc & Để trống \texttt{ho\_ten} & Form chặn tại frontend & Đạt \\ \hline
4 & Sửa thông tin & Đổi \texttt{cap\_bac} & Lưu thành công & Đạt \\ \hline
5 & Chuyển đơn vị & Chọn đơn vị mới & Lịch sử khen thưởng giữ nguyên, bộ đếm cập nhật & Đạt \\ \hline
\end{tabular}
\end{table}

\begin{table}[H]
\centering
\caption{Kiểm thử chức năng Đề xuất khen thưởng}
\label{tab:kt-dexuat}
\begin{tabular}{|c|p{3.8cm}|p{3.3cm}|p{3.8cm}|c|}
\hline
\textbf{STT} & \textbf{Chức năng} & \textbf{Đầu vào} & \textbf{Đầu ra mong muốn} & \textbf{Kết quả} \\ \hline
1 & Tạo đề xuất hợp lệ & 5 quân nhân, CSTĐCS & Đề xuất ở trạng thái PENDING & Đạt \\ \hline
2 & Manager ngoài đơn vị & Quân nhân đơn vị khác & Backend trả 403 & Đạt \\ \hline
3 & Quân nhân không đủ điều kiện & 1 năm CSTĐCS đề nghị BKBTBQP & Cảnh báo trước khi gửi & Đạt \\ \hline
4 & Lưu nháp khi tải lại trang & Reload trang giữa chừng & Dữ liệu giữ nguyên (loại, quân nhân, danh hiệu, ghi chú); xoá tự động sau khi gửi hoặc đóng tab & Đạt \\ \hline
\end{tabular}
\end{table}

\begin{table}[H]
\centering
\caption{Kiểm thử chức năng Phê duyệt đề xuất}
\label{tab:kt-pheduyet}
\begin{tabular}{|c|p{3.8cm}|p{3.3cm}|p{3.8cm}|c|}
\hline
\textbf{STT} & \textbf{Chức năng} & \textbf{Đầu vào} & \textbf{Đầu ra mong muốn} & \textbf{Kết quả} \\ \hline
1 & Phê duyệt thông thường & Đề xuất hợp lệ & Trạng thái APPROVED, ghi vào bảng tương ứng & Đạt \\ \hline
2 & Phê duyệt kèm gắn quyết định & Số quyết định + tháng -- năm nhận & Đề xuất gắn quyết định, FE quân nhân hiển thị đúng dữ liệu & Đạt \\ \hline
3 & Từ chối đề xuất & Bấm ``Từ chối'' kèm lý do & REJECTED, gửi thông báo Manager & Đạt \\ \hline
4 & Vi phạm quy tắc BKTTCP một lần duy nhất & Đề nghị BKTTCP cho người đã có & Giao dịch huỷ bỏ & Đạt \\ \hline
5 & Hai Admin cùng phê duyệt & Hai phiên đồng thời & Chỉ một thành công, người sau nhận 409 & Đạt \\ \hline
\end{tabular}
\end{table}

\begin{table}[H]
\centering
\caption{Kiểm thử chức năng Tính lại điều kiện khen thưởng hằng năm}
\label{tab:kt-tinhlai}
\begin{tabular}{|c|p{3.8cm}|p{3.3cm}|p{3.8cm}|c|}
\hline
\textbf{STT} & \textbf{Chức năng} & \textbf{Đầu vào} & \textbf{Đầu ra mong muốn} & \textbf{Kết quả} \\ \hline
1 & Quân nhân chưa nhận & 2 năm CSTĐCS liên tục & Đề nghị BKBTBQP, gợi ý giải thích & Đạt \\ \hline
2 & Đã nhận BKTTCP (một lần duy nhất) & Đã có BKTTCP & Thông điệp ``Đã có BKTTCP. Phần mềm chưa hỗ trợ...'' & Đạt \\ \hline
3 & Bỏ lỡ đợt & 5 năm CSTĐCS từ BKBTBQP cuối & \texttt{missedBkbqp = 2}, vẫn đề nghị BKBTBQP mới & Đạt \\ \hline
4 & Khung 3 năm CSTĐTQ & BKBTBQP cách 4 năm & BKBTBQP đó không tính, yêu cầu mới & Đạt \\ \hline
5 & Thiếu NCKH năm thứ 3 & 3 năm CSTĐCS, thiếu NCKH năm 2 & Không đủ, gợi ý cần NCKH & Đạt \\ \hline
\end{tabular}
\end{table}

\begin{table}[H]
\centering
\caption{Kiểm thử chức năng Nhập Excel hàng loạt cho danh hiệu hằng năm}
\label{tab:kt-excel}
\begin{tabular}{|c|p{3.8cm}|p{3.3cm}|p{3.8cm}|c|}
\hline
\textbf{STT} & \textbf{Chức năng} & \textbf{Đầu vào} & \textbf{Đầu ra mong muốn} & \textbf{Kết quả} \\ \hline
1 & Nhập 50 bản ghi CSTĐCS hợp lệ & Tệp \texttt{.xlsx} đúng mẫu & 50 bản ghi được tạo, ghi nhật ký hệ thống & Đạt \\ \hline
2 & Có dòng trùng (cùng quân nhân và năm) & 1/50 dòng trùng khoá & Xem trước hiển thị 49 valid + 1 error & Đạt \\ \hline
3 & Confirm khi DB thay đổi giữa hai bước & Dữ liệu chen vào sau preview & Giao dịch huỷ bỏ, không bản ghi nào được ghi & Đạt \\ \hline
4 & Tệp sai định dạng & Tệp \texttt{.doc} thay vì \texttt{.xlsx} & Trả mã 400 kèm mô tả lỗi & Đạt \\ \hline
\end{tabular}
\end{table}

\begin{table}[H]
\centering
\caption{Kiểm thử chức năng Quản lý đơn vị và chức vụ}
\label{tab:kt-donvi}
\begin{tabular}{|c|p{3.8cm}|p{3.3cm}|p{3.8cm}|c|}
\hline
\textbf{STT} & \textbf{Chức năng} & \textbf{Đầu vào} & \textbf{Đầu ra mong muốn} & \textbf{Kết quả} \\ \hline
1 & Thêm cơ quan đơn vị mới & Tên + mã & Có trong cây & Đạt \\ \hline
2 & Xoá đơn vị trực thuộc có quân nhân & Đơn vị trực thuộc vẫn có người & Backend chặn xoá, trả 409 & Đạt \\ \hline
3 & Sửa hệ số chức vụ & Đổi 0{,}8 $\rightarrow$ 0{,}9 & Lưu thành công, hồ sơ HCBVTQ tự cập nhật & Đạt \\ \hline
\end{tabular}
\end{table}

\begin{table}[H]
\centering
\caption{Kiểm thử chức năng Quản trị (sao lưu, nhật ký)}
\label{tab:kt-quantri}
\begin{tabular}{|c|p{3.8cm}|p{3.3cm}|p{3.8cm}|c|}
\hline
\textbf{STT} & \textbf{Chức năng} & \textbf{Đầu vào} & \textbf{Đầu ra mong muốn} & \textbf{Kết quả} \\ \hline
1 & Sao lưu thủ công & Bấm ``Sao lưu ngay'' & Tệp .sql mới trong \texttt{backups/} & Đạt \\ \hline
2 & Manager xem nhật ký backup & Mở trang nhật ký & Không thấy mục \texttt{resource = backup} & Đạt \\ \hline
3 & Đổi lịch sao lưu & 24h $\rightarrow$ 12h & Backup tự chạy theo lịch mới & Đạt \\ \hline
4 & Khôi phục từ tệp sao lưu & Chọn tệp + xác nhận & DB trở về trạng thái lưu trữ & Đạt \\ \hline
\end{tabular}
\end{table}

\subsection{Kiểm thử tương thích}
\label{subsection:4.4.2}

Hệ thống được kiểm thử trên năm dòng máy phổ biến với các tổ hợp hệ điều hành và trình duyệt khác nhau nhằm đánh giá khả năng hiển thị và vận hành đồng nhất.

\begin{table}[H]
\centering
\caption{Kiểm thử tương thích các dòng máy}
\label{tab:kt-trinhduyet}
\renewcommand{\arraystretch}{1.3}
\begin{tabular}{|p{3.2cm}|p{4.2cm}|p{3.6cm}|p{2.4cm}|}
\hline
\textbf{Thiết bị} & \textbf{Phần cứng} & \textbf{Phần mềm} & \textbf{Kết quả} \\ \hline
MacBook Pro M1 Pro 2021 & Apple M1 Pro, RAM 16 GB, màn hình 14,2 inch Liquid Retina XDR & macOS Sonoma 14, Safari 17 & Đạt (giới hạn nhỏ ở DatePicker) \\ \hline
Lenovo Legion 5 2022 & AMD Ryzen 7 5800H, RAM 16 GB, màn hình 15,6 inch FHD 165 Hz & Windows 11, Microsoft Edge 119 & Đạt \\ \hline
Lenovo IdeaPad Gaming 3 & AMD Ryzen 5 5600H, RAM 8 GB, màn hình 15,6 inch FHD & Windows 11, Google Chrome 120 & Đạt \\ \hline
Dell Inspiron 15 3520 & Intel Core i5-1235U, RAM 16 GB, màn hình 15,6 inch FHD & Windows 10, Mozilla Firefox 121 & Đạt \\ \hline
Mac mini M4 2024 & Apple M4, RAM 16 GB, màn hình ngoài 23,8 inch FHD & macOS Sequoia 15, Google Chrome 120 & Đạt \\ \hline
\end{tabular}
\end{table}

\section{Triển khai}
\label{section:4.5}

\subsection{Yêu cầu phần cứng và phần mềm}
\label{subsection:4.5.1}

Hệ thống có thể triển khai trên một máy chủ Linux duy nhất với cấu hình tối thiểu sau:

\begin{itemize}
    \item CPU: 2 nhân, 2{,}0 GHz trở lên.
    \item RAM: 4 GB (khuyến nghị 8 GB cho môi trường vài chục người dùng).
    \item Ổ cứng: 50 GB SSD cho dữ liệu và bản sao lưu.
    \item Hệ điều hành: Ubuntu Server 22.04 LTS hoặc Debian 12.
    \item Phần mềm phụ thuộc: Node.js 20 LTS, PostgreSQL 15.
\end{itemize}

\subsection{Hướng dẫn triển khai môi trường phát triển}
\label{subsection:4.5.2}

Các bước cài đặt trên máy phát triển cá nhân được trình bày tuần tự dưới đây.

\textbf{Bước 1.} Cài đặt phụ thuộc Node bằng lệnh \texttt{npm install} lần lượt trong hai thư mục \texttt{BE-QLKT} và \texttt{FE-QLKT}. Lệnh này tải các gói phụ thuộc liệt kê tại tệp \texttt{package.json} của mỗi module.

\textbf{Bước 2.} Sao chép tệp cấu hình mẫu \texttt{.env.example} thành \texttt{.env} trong cả hai thư mục, rồi chỉnh các giá trị thiết yếu: \texttt{DATABASE\_URL} cho chuỗi kết nối PostgreSQL, \texttt{JWT\_SECRET} cho khoá phát hành Access Token, \texttt{JWT\_REFRESH\_SECRET} cho khoá phát hành Refresh Token, \texttt{CLIENT\_URL} cho địa chỉ frontend.

\textbf{Bước 3.} Khởi tạo cơ sở dữ liệu trong thư mục \texttt{BE-QLKT} bằng lệnh \texttt{npx prisma migrate dev} để áp dụng các tệp migration đã có; tiếp theo chạy lệnh \texttt{npm run init-super-admin} để tạo tài khoản SuperAdmin mặc định.

\textbf{Bước 4.} Khởi chạy môi trường phát triển bằng \texttt{npm run dev} ở \texttt{BE-QLKT} cho backend tại cổng 4000 và lệnh tương tự ở \texttt{FE-QLKT} cho frontend tại cổng 3000.

Sau bước 4, mở trình duyệt tới \url{http://localhost:3000} và đăng nhập bằng tài khoản SuperAdmin mặc định để bắt đầu cấu hình các đơn vị, chức vụ và tài khoản phụ trợ.

\subsection{Triển khai sản xuất với PM2}
\label{subsection:4.5.3}

Trên máy chủ sản xuất, sản phẩm được vận hành bằng PM2, một công cụ quản lý tiến trình hỗ trợ tự khởi động lại khi tiến trình bị dừng đột ngột, ghi log có thời gian và theo dõi tài nguyên hệ thống. Tệp \texttt{ecosystem.config.js} ở thư mục gốc của dự án khai báo hai tiến trình ứng dụng: \texttt{be-qlkt} chạy bản đã build (\texttt{dist/index.js}) trong thư mục \texttt{BE-QLKT} ở cổng 4000 và \texttt{fe-qlkt} chạy \texttt{next start} trong thư mục \texttt{FE-QLKT} ở cổng 3000. Cả hai tiến trình được đặt biến môi trường \texttt{NODE\_ENV} là \texttt{production}.

Khởi chạy bằng lệnh \texttt{pm2 start ecosystem.config.js}, sau đó dùng \texttt{pm2 save} để ghi nhớ danh sách tiến trình đang chạy và \texttt{pm2 startup} để cấu hình tự khởi động khi máy chủ khởi động lại. Người dùng trong mạng nội bộ truy cập giao diện do tiến trình frontend phục vụ ở cổng 3000; trình duyệt gọi trực tiếp tới backend ở cổng 4000 qua REST API và Socket.IO.

\subsection{Sao lưu và giám sát}
\label{subsection:4.5.4}

Tác vụ sao lưu chạy bên trong tiến trình backend qua \texttt{node-cron}, lịch và trạng thái cấu hình ngay trong DevZone. Lịch mặc định là \texttt{0 1 1 * *} (01:00 ngày 1 hằng tháng); cho mỗi lần chạy, dịch vụ đọc 21 bảng dữ liệu, ghép chuỗi câu lệnh \texttt{INSERT INTO \dots} đã escape rồi ghi tệp vào \texttt{backups/YYYY-MM-DD\_HH-mm-ss.sql}. Các tệp cũ vượt \texttt{backup\_retention\_days} (mặc định 15) được tự động xoá ở bước cleanup. PM2 cung cấp \texttt{pm2 logs} để xem log thời gian thực; với giám sát dài hạn, log gom vào \texttt{/var/log/qlkt/} và có thể tích hợp với hệ thống giám sát tập trung của Học viện trong tương lai.


\end{document}
